\documentclass[11pt,openany]{book}

% GraphHunter — Libro «Diseño completo»
% Pedagogical design doc: complete walkthrough of GraphHunter's
% architecture, algorithms, and forward-looking stages.
%
% Accent color: morado oscuro (Diseño) for visual distinction
% from the other four books.

\usepackage{fontspec}
\setmainfont{Latin Modern Roman}
\setsansfont{Latin Modern Sans}
\setmonofont{Latin Modern Mono}
\usepackage[spanish,es-tabla,es-lcroman,es-noshorthands]{babel}
\usepackage{csquotes}

\usepackage[a4paper,margin=2.4cm]{geometry}

\usepackage{amsmath, amssymb, amsthm}
\usepackage{mathtools}

\usepackage[ruled,linesnumbered,algochapter]{algorithm2e}
\SetKwComment{Comment}{$\triangleright$ }{}
\SetKwFor{For}{for}{do}{end}
\SetKwFor{While}{while}{do}{end}
\SetKwIF{If}{ElseIf}{Else}{if}{then}{else if}{else}{end}
\SetKwRepeat{Repeat}{repeat}{until}

\usepackage{listings}
\usepackage{xcolor}
\definecolor{accent}{HTML}{4B2E83}       % morado oscuro (Diseño)
\definecolor{accentlight}{HTML}{EAE0F4}
\definecolor{codecomment}{rgb}{0.35,0.45,0.45}
\definecolor{codekw}{rgb}{0.55,0.10,0.55}
\definecolor{codestr}{rgb}{0.10,0.45,0.10}
\lstdefinelanguage{Rust}{
    morekeywords={fn,let,mut,pub,struct,enum,impl,trait,where,if,else,for,while,
                  loop,match,return,use,mod,const,static,self,Self,as,in,ref,move,
                  async,await,dyn,unsafe,crate,super,extern,type,break,continue},
    morecomment=[l]{//},
    morecomment=[s]{/*}{*/},
    morestring=[b]",
    morestring=[b]',
}
\lstset{
    basicstyle=\ttfamily\small,
    keywordstyle=\color{codekw}\bfseries,
    commentstyle=\color{codecomment}\itshape,
    stringstyle=\color{codestr},
    showstringspaces=false,
    frame=single,
    numbers=none,
    captionpos=b,
    breaklines=true,
    tabsize=4,
    % Listings comen Unicode salvo que les digas literate→glyph.
    % Mapeamos los símbolos griegos / matemáticos que aparecen en
    % comentarios de Rust dentro de los listings.
    extendedchars=true,
    literate=%
        {α}{{$\alpha$}}1
        {β}{{$\beta$}}1
        {ρ}{{$\rho$}}1
        {Σ}{{$\Sigma$}}1
        {≤}{{$\leq$}}1
        {≥}{{$\geq$}}1
        {→}{{$\rightarrow$}}1
        {↔}{{$\leftrightarrow$}}1,
}

\usepackage{hyperref}
\hypersetup{
    colorlinks=true,
    linkcolor=accent,
    citecolor=black,
    urlcolor=accent,
    pdftitle={GraphHunter — Diseño completo},
    pdfauthor={Lucas Sotomayor / Base4Security},
    bookmarks=true,
    bookmarksopen=true,
    bookmarksnumbered=true,
}
\usepackage{bookmark}

\usepackage{booktabs}
\usepackage{longtable}
\usepackage{float}
\usepackage{tikz}
\usetikzlibrary{arrows.meta, positioning, shapes.geometric, fit, backgrounds, calc, shapes.misc, decorations.pathreplacing}

% Box style for "intuición" callouts — used to give a plain-language
% summary before each formal definition.
\usepackage[most]{tcolorbox}
\newtcolorbox{intuicion}{
    colback=accentlight,
    colframe=accent,
    arc=2mm,
    boxrule=0.6pt,
    left=4pt, right=4pt, top=4pt, bottom=4pt,
    fonttitle=\bfseries,
    title=Intuición
}

\theoremstyle{plain}
\newtheorem{theorem}{Teorema}[chapter]
\newtheorem{lemma}[theorem]{Lema}

\theoremstyle{definition}
\newtheorem{definition}[theorem]{Definición}

\theoremstyle{remark}
\newtheorem*{invariant}{Invariante}
\newtheorem*{nota}{Nota}

\newcommand{\filepath}[1]{\texttt{\small\detokenize{#1}}}
\newcommand{\code}[1]{\texttt{\detokenize{#1}}}

\newenvironment{enelcodigo}{%
    \par\vspace{0.4em}%
    \noindent\textbf{En el código.}\nopagebreak\par\nopagebreak\vspace{0.2em}%
}{\par\vspace{0.4em}}

\usepackage{titlesec}
\titleformat{\chapter}[display]
  {\normalfont\huge\bfseries\color{accent}}
  {\chaptertitlename~\thechapter}
  {20pt}
  {\Huge}

\newcommand{\TODO}[1]{\textbf{\textcolor{red}{[TODO: #1]}}}
\newcommand{\passthrough}[1]{#1}


\title{\Huge\bfseries Graph Hunter\\[0.4em]\color{accent}\Large Diseño completo}
\author{Lucas Sotomayor\\Base4Security}
\date{Abril 2026}

\begin{document}

\frontmatter
\maketitle

\begin{flushleft}
\textbf{¿Qué es este libro?}\\
Este libro es la guía pedagógica del diseño completo de Graph Hunter,
desde el modelo de datos hasta la pila de optimizaciones de \emph{matching}
de las tiers \emph{S}, \emph{M} y \emph{L} entregadas en abril 2026.
A diferencia del libro \emph{Arquitectura} (que documenta la estructura
de \emph{crates}), del libro \emph{Especificación} (que enumera \emph{wire
contracts}), y del libro \emph{Manual} (que enseña al analista a usarlo),
este libro toma decisiones algorítmicas y las explica con \emph{intuición},
\emph{diagramas} y \emph{pseudocódigo} antes de entrar al código.
\end{flushleft}

\tableofcontents

\mainmatter

\chapter{¿Qué es Graph Hunter?}
\label{cap:intro}

\begin{intuicion}
Graph Hunter es un motor de \emph{threat hunting} basado en grafos
temporales. Convierte \emph{logs} de seguridad (Sysmon, Sentinel,
EVTX, Fortinet) en un grafo dirigido con tiempo, donde los \emph{nodos}
son entidades (procesos, IPs, usuarios, archivos) y las \emph{aristas}
son eventos temporales que las relacionan. Sobre ese grafo el analista
expresa hipótesis (``buscame todos los \emph{User} que se autenticaron
en un \emph{Host} que después se conectó a una \emph{IP} sospechosa
en menos de 30~minutos'') en una DSL que el motor ejecuta como
\emph{subgraph matching} con restricciones temporales.
\end{intuicion}

\section{Por qué un grafo y no un \emph{SIEM} clásico}

Los \emph{SIEM} tradicionales son motores de búsqueda full-text sobre
\emph{logs}: \texttt{username=admin AND host=db-01}. Funcionan cuando
el analista sabe \emph{qué busca}. La debilidad aparece cuando la
señal es \emph{multi-hop}: un usuario inocuo, un host comprometido,
y una conexión saliente sospechosa. Cada uno por separado pasa los
filtros; sólo la \emph{cadena} entre ellos delata el ataque.

Modelar el grafo explícitamente permite expresar y matchear esas
cadenas en una sola consulta. Donde el SIEM hace \texttt{n}
\emph{joins} sobre \texttt{n} tablas — cada uno con su propio índice
— el grafo recorre aristas en su \emph{adjacency list} y aplica
filtros temporales \emph{in-line}.

\begin{figure}[H]
\centering
\begin{tikzpicture}[
    node distance=1.4cm,
    every node/.style={font=\small},
    entity/.style={draw=accent, thick, circle, minimum size=8mm,
                   inner sep=1pt, fill=accentlight},
    edge/.style={->, >=Latex, accent, thick}
]
\node[entity] (u) {\strut $u$};
\node[entity, right=2.4cm of u] (h) {\strut $h$};
\node[entity, right=2.4cm of h] (ip) {\strut $\mathit{ip}$};

\draw[edge] (u) -- node[above]{\footnotesize\texttt{Auth}\,@\,$t_1$} (h);
\draw[edge] (h) -- node[above]{\footnotesize\texttt{Connect}\,@\,$t_2$} (ip);

\node[draw=accent, dashed, fit=(u)(h)(ip), inner sep=4mm,
      label={[font=\footnotesize\itshape, accent]below:%
             hipótesis: \texttt{User -[Auth]-> Host -[Connect]-> IP \{within=30m\}}}]{};
\end{tikzpicture}
\caption{Ejemplo mínimo: una hipótesis de dos saltos. El motor enumera
todas las ternas $(u, h, \mathit{ip})$ del grafo donde existen las dos
aristas con $t_1 < t_2$ y $t_2 - t_1 \leq 30$~minutos.}
\label{fig:intro-2hop}
\end{figure}

\section{Pipeline en cuatro pasos}

Graph Hunter está dividido en cuatro etapas que el analista atraviesa
linealmente:

\begin{enumerate}
    \item \textbf{Ingesta}. Una fuente (\emph{Sysmon}, \emph{Sentinel},
    \emph{EVTX}, CSV genérico) emite eventos. El parser produce
    \emph{triples} \texttt{(source, relation, dest, timestamp,
    metadata)} que el \emph{ingestor} consolida en el grafo.
    \item \textbf{Hipótesis}. El analista escribe una expresión DSL.
    El \emph{parser} la convierte en una \texttt{Hypothesis}
    estructurada (cadena de \texttt{HypothesisStep}, predicados,
    bindings, \emph{shape}).
    \item \textbf{Matching}. El \emph{planner} elige el operador
    físico adecuado (DFS legacy, LFTJ-K4, Yannakakis chain,
    Yannakakis tree-walk) según la forma de la hipótesis y \emph{flags}
    de \emph{env}. El operador devuelve \texttt{Vec<HuntResult>}.
    \item \textbf{Scoring + paginación}. El \emph{anomaly scorer}
    rankea los \emph{paths} por rareza; el \emph{frontend} paginación
    sobre la \emph{cache} de resultados.
\end{enumerate}

\begin{figure}[H]
\centering
\begin{tikzpicture}[
    node distance=0.8cm and 0.9cm,
    every node/.style={font=\footnotesize, align=center},
    stage/.style={draw=accent, thick, rounded corners=2mm,
                  minimum width=2.6cm, minimum height=1.0cm,
                  fill=accentlight},
    arr/.style={->, >=Latex, accent, thick}
]
\node[stage] (ingest) {Ingesta\\\textit{Sysmon, EVTX, CSV}};
\node[stage, right=of ingest] (graph) {Grafo temporal\\\textit{en RAM + spill}};
\node[stage, right=of graph] (parse) {Parser DSL\\\textit{Hypothesis AST}};
\node[stage, right=of parse] (match) {Matcher\\\textit{Yannakakis $\bigm|$ DFS $\bigm|$ LFTJ}};
\node[stage, right=of match] (score) {Scoring +\\paginación};

\draw[arr] (ingest) -- (graph);
\draw[arr] (graph) -- (parse);
\draw[arr] (parse) -- (match);
\draw[arr] (match) -- (score);

\node[below=0.3cm of ingest, font=\footnotesize\itshape, accent]
    {ev/sec $\Rightarrow$ aristas/sec};
\node[below=0.3cm of graph, font=\footnotesize\itshape, accent]
    {28~GB EVTX viable};
\node[below=0.3cm of match, font=\footnotesize\itshape, accent]
    {3-hop @ $|V|$=10K en 10~ms};
\end{tikzpicture}
\caption{Pipeline alto-nivel. Cada etapa tiene su propia hoja de ruta
de optimizaciones; las tiers \emph{S} (\emph{small}), \emph{M}
(\emph{medium}) y \emph{L} (\emph{large}) que cubre este libro
afectan sobre todo a \emph{Matcher} y \emph{Grafo temporal}.}
\label{fig:intro-pipeline}
\end{figure}

\section{Tesis transversal: \emph{plug-and-play} como promesa verificable}
\label{sec:intro-plug-and-play}

Antes de entrar al stack, conviene hacer explícita la \emph{tesis
transversal} que ata todas las decisiones de diseño que vienen a
continuación: \textbf{ninguna optimización entregada requiere que
el usuario cambie su \emph{config}}. Cada \emph{flag} S/M/L que
ships se setea \emph{default-on} al \emph{startup} del \emph{binary},
cada \emph{cache} es transparente, cada índice se construye solo,
cada índice grande se \emph{spills} a disco solo. El usuario abre
la app, ve la pildora ``Perf: M+L1'' verde, y todo funciona.

Esta restricción no es cosmética: es lo que hace que el motor sea
auditable. Si después de un \emph{ship} de optimización el analista
tuvo que setear \texttt{GRAPHHUNTER\_LFTJ\_PLANNER=1} para verla,
la promesa se rompió --- la optimización no está realmente
entregada, está \emph{disponible}. El capítulo~\ref{cap:plug-and-play}
codifica la regla en cinco pasos concretos (default-on al startup,
\emph{surface} en UI, \emph{soak window} para los \emph{flags}
peligrosos, doc obliga, \emph{build capabilities} reflejadas en
\texttt{cmd\_get\_perf\_status}). Cuando leas en los capítulos que
siguen frases como ``X dispatcha automáticamente'' o ``Y se
construye \emph{lazy} en el primer hunt'', son aplicaciones de
esta misma regla.

\section{Lo que cubre este libro}

\begin{itemize}
    \item Capítulo~\ref{cap:stack}: el \emph{stack} de \emph{crates}.
    Quién depende de quién y por qué, con el diagrama de capas.
    \item Capítulo~\ref{cap:datos}: el modelo de datos
    (\texttt{Entity}, \texttt{Relation}, \texttt{CompactRelation},
    el \emph{interner} de \texttt{StrId}, \texttt{HuntResult}).
    \item Capítulo~\ref{cap:ingest}: el pipeline de ingesta y
    finalización (\texttt{sort\_edges\_by\_timestamp}, \emph{spill},
    \emph{streaming backend}).
    \item Capítulo~\ref{cap:matcher}: el camino del \emph{matching}
    operador-por-operador. DFS, LFTJ-K4, Yannakakis chain,
    Yannakakis tree-walk. Cada uno con su algoritmo, sus
    invariantes y su \emph{when-to-dispatch}.
    \item Capítulo~\ref{cap:indexes}: NLF y k-hop reachability,
    incluyendo el \emph{lazy build} y la \emph{sparse mmap-spill}.
    \item Capítulo~\ref{cap:dsl}: la DSL, incluido el separador
    \texttt{;} para \emph{shapes} no-cadena de L1 stage 3c.
    \item Capítulo~\ref{cap:l3-future}: diseño de las stages
    pendientes \emph{2b} (delta khop\_reach via Kankanamge) y
    \emph{2c} (streaming como \emph{canonical store}).
    \item Capítulo~\ref{cap:plug-and-play}: la promesa
    \emph{plug-and-play} y cómo se materializa en el frontend.
\end{itemize}

\chapter{Stack arquitectónico}
\label{cap:stack}

\begin{intuicion}
El repositorio se divide en \emph{layers} concéntricos. La regla de
oro: cada \emph{layer} sólo puede importar capas iguales o más
internas. \texttt{core} no toca \texttt{platform}; \texttt{platform}
no toca \texttt{apps}. Esa restricción se valida con
\filepath{scripts/check-layering.sh} en cada CI run.
\end{intuicion}

\section{Diagrama de capas}

\begin{figure}[H]
\centering
\begin{tikzpicture}[
    node distance=0.6cm,
    every node/.style={font=\footnotesize, align=center},
    layer/.style={draw=accent, thick, rounded corners=2mm,
                  minimum height=1cm, fill=accentlight},
    crate/.style={draw=accent!70, thick, rounded corners=1mm,
                  minimum height=0.7cm, font=\scriptsize\ttfamily,
                  fill=white}
]
% Apps layer
\node[layer, minimum width=14cm] (apps) {\textbf{apps/}\,---\,binarios end-user};
\node[crate, below right=0.3cm and -7.5cm of apps.south west,
      minimum width=2.5cm] (cli)   {cli};
\node[crate, right=0.4cm of cli, minimum width=2.5cm] (tauri) {tauri};
\node[crate, right=0.4cm of tauri, minimum width=2.5cm] (gateway) {go-gateway};
\node[crate, right=0.4cm of gateway, minimum width=2.5cm] (mcp)   {mcp};

% Platform layer
\node[layer, minimum width=14cm, below=2.0cm of apps] (plat)
    {\textbf{platform/}\,---\,API canónica + parsers + sources};
\node[crate, below right=0.3cm and -7.5cm of plat.south west,
      minimum width=2.0cm] (api)      {api};
\node[crate, right=0.3cm of api, minimum width=1.8cm] (dsl)      {dsl};
\node[crate, right=0.3cm of dsl, minimum width=2.0cm] (sources)  {sources};
\node[crate, right=0.3cm of sources, minimum width=1.8cm] (parsers) {parsers};
\node[crate, right=0.3cm of parsers, minimum width=2.0cm] (canonical) {canonical};
\node[crate, right=0.3cm of canonical, minimum width=1.6cm] (vrl)  {vrl};

% Core layer
\node[layer, minimum width=14cm, below=2.0cm of plat] (core)
    {\textbf{core/}\,---\,grafo temporal + matcher + scoring};
\node[crate, below right=0.3cm and -6.5cm of core.south west,
      minimum width=3.4cm] (engine)  {graph-engine};
\node[crate, right=0.4cm of engine, minimum width=3.0cm] (ffi)     {matcher-ffi};
\node[crate, right=0.4cm of ffi, minimum width=2.4cm] (gnn)        {gnn};

% Arrows between layers (apps -> platform -> core)
\draw[->, >=Latex, accent, very thick] (apps.south) -- (plat.north);
\draw[->, >=Latex, accent, very thick] (plat.south) -- (core.north);
\end{tikzpicture}
\caption{Layering. Las flechas representan la dirección permitida de
\emph{imports}; \texttt{apps} puede usar \texttt{platform} y
\texttt{core}, \texttt{platform} sólo \texttt{core}, \texttt{core} no
ve hacia arriba. Excepciones puntuales se documentan en
\filepath{docs/arquitectura/layering-exceptions.txt} y obligan ADR.}
\label{fig:stack-layers}
\end{figure}

\section{Por qué este orden}

\begin{itemize}
    \item \textbf{core} es la pieza algorítmica pura. No conoce a
    Tauri, ni a Sentinel, ni a HTTP. Eso le permite testarse en
    aislamiento (\emph{property tests}, benches \emph{criterion}) y
    reusarse desde múltiples \emph{transports} sin acoplamiento.
    \item \textbf{platform} agrupa los \emph{contracts} compartidos
    entre transports: la \emph{canonical API}
    (\texttt{GraphHunterApi}), el parser de la DSL, el adapter de
    \emph{sources} (Sentinel, Defender XDR, Data Lake), los parsers
    concretos (Cognito, FortiAnalyzer, IIS), y la \emph{canonical
    pipeline} (VRL + provenance).
    \item \textbf{apps} son los binarios usables: el \emph{daemon}
    Tauri (con su frontend React+Vite), el CLI que el \emph{Go
    gateway} invoca como sub-process, y el plugin MCP. Todos
    consumen \texttt{Arc<GraphHunterApi>} de \texttt{platform/api}.
\end{itemize}

\section{Cómo se decide en qué \emph{crate} va algo}

Tres preguntas guían la decisión:

\begin{enumerate}
    \item \emph{¿Lo necesita más de un transport?} Sí $\Rightarrow$
    \texttt{platform}. No $\Rightarrow$ \texttt{apps}.
    \item \emph{¿Tiene contrato externo (HTTP, MCP, FFI)?} Sí
    $\Rightarrow$ siempre \texttt{platform}, nunca dentro de un
    \texttt{apps/} específico (porque cualquier \emph{transport}
    puede querer reusarlo).
    \item \emph{¿Es algorítmico puro y libre de I/O o de
    \emph{transport}?} Sí $\Rightarrow$ \texttt{core}.
\end{enumerate}

\section{Flujo de un \emph{hunt} a través del \emph{stack}}

\begin{figure}[H]
\centering
\begin{tikzpicture}[
    node distance=0.5cm,
    every node/.style={font=\scriptsize, align=center},
    box/.style={draw=accent, thick, rounded corners=1.5mm,
                minimum height=0.7cm, minimum width=2.6cm,
                fill=accentlight}
]
\node[box] (frontend) {React frontend\\\itshape App.tsx};
\node[box, below=of frontend] (cmd)      {Tauri command\\\itshape cmd\_run\_hunt};
\node[box, below=of cmd] (api)           {GraphHunterApi\\\itshape run\_hunt};
\node[box, below=of api] (graph)         {GraphHunter\\\itshape search\_temporal\_pattern};
\node[box, below=of graph] (planner)     {planner\\\itshape DispatchPlan};
\node[box, below left=of planner, xshift=-0.6cm] (yann) {yannakakis};
\node[box, below=of planner] (dfs)       {DFS legacy};
\node[box, below right=of planner, xshift=0.6cm] (lftj) {LFTJ-K4};

\draw[->, >=Latex, accent, thick] (frontend) -- (cmd);
\draw[->, >=Latex, accent, thick] (cmd) -- (api);
\draw[->, >=Latex, accent, thick] (api) -- (graph);
\draw[->, >=Latex, accent, thick] (graph) -- (planner);
\draw[->, >=Latex, accent, thick] (planner) -- (yann);
\draw[->, >=Latex, accent, thick] (planner) -- (dfs);
\draw[->, >=Latex, accent, thick] (planner) -- (lftj);
\end{tikzpicture}
\caption{Camino de un \emph{hunt} desde el botón ``Run'' del frontend
hasta el operador físico que el \emph{planner} elija. Las tres
ramas inferiores son los operadores existentes; el \emph{planner} es
el árbitro.}
\label{fig:stack-hunt-flow}
\end{figure}

\chapter{Modelo de datos}
\label{cap:datos}

\begin{intuicion}
El grafo es un \emph{multigrafo dirigido temporal} sobre entidades.
``Multi'' = puede haber varias aristas entre los mismos dos nodos
(por ejemplo, el mismo proceso lanzando otro varias veces). ``Dirigido''
= cada arista tiene origen y destino fijos, no son simétricas.
``Temporal'' = cada arista lleva un \texttt{timestamp} y los algoritmos
explotan ese orden. ``Sobre entidades'' = los nodos no son strings sueltos
sino tuplas tipadas \texttt{(EntityType, identifier)}.
\end{intuicion}

\section{\texttt{Entity}, \texttt{Relation}, \texttt{StrId}}

\begin{lstlisting}[language=Rust,caption={Tipos públicos del modelo. Los campos opcionales se omiten en el ejemplo para foco.}]
pub struct Entity {
    pub id: String,                 // ej. "10.0.0.42" o "lsass.exe"
    pub entity_type: EntityType,    // ej. IP, Process, User
    pub metadata: HashMap<...>,     // campos opcionales del log
}

pub struct Relation {
    pub source_id: String,
    pub dest_id: String,
    pub rel_type: RelationType,     // ej. Auth, Connect, Spawn
    pub timestamp: i64,             // segundos epoch
    pub metadata: HashMap<...>,
}
\end{lstlisting}

\texttt{Entity::id} y \texttt{Relation::source\_id}/\texttt{dest\_id}
son \texttt{String}s humanos. Internamente, el motor los pasa por un
\emph{interner} (\texttt{StringInterner}) que devuelve un
\texttt{StrId} (\texttt{u32}) compacto. La \texttt{StringInterner}
guarda cada string una sola vez en un \texttt{Vec<Arc<str>>};
resoluciones futuras devuelven \texttt{Arc<str>} sin reasignación.

\begin{intuicion}
Pensá el \emph{interner} como una tabla de hash bidireccional con
identidad estable. Cada string nuevo entra con un \texttt{StrId}
secuencial; el motor pasa \texttt{StrId}s en lugar de \texttt{String}s
durante todo el \emph{matching} (4~bytes vs.\ ~24~bytes/punter +
\emph{heap} en cada copia). Sólo al \emph{emit} de un resultado se
resuelven a \texttt{Arc<str>}.
\end{intuicion}

\section{\texttt{CompactRelation}: 32~bytes/arista}

El \texttt{Relation} público (humano) pesa $\sim 200$~bytes en RAM.
Para soportar EVTX de 28~GB ($\sim 100$M aristas) sin reventar la
\emph{heap}, el motor usa un \texttt{CompactRelation} de 32~bytes:

\begin{lstlisting}[language=Rust]
#[repr(C)]
pub struct CompactRelation {
    pub timestamp: i64,           // 8 B
    pub metadata_offset: u64,     // 8 B (offset al MetadataStore)
    pub source_sid: StrId,        // 4 B
    pub dest_sid: StrId,          // 4 B
    pub ingested_at_delta: u32,   // 4 B (segundos desde INGEST_EPOCH_BASE)
    pub dataset_tag: u16,         // 2 B
    pub rel_type_tag: u8,         // 1 B
    // 1 B padding implícito
}
\end{lstlisting}

\texttt{metadata\_offset} es un \emph{handle} al
\texttt{MetadataStore} \emph{append-only}; \texttt{0} significa
``sin metadata'' y se chequea sin \emph{lookup}. El \texttt{u8} de
\texttt{rel\_type\_tag} permite el \emph{matching} sin resolver el
\texttt{enum} entero.

\begin{figure}[H]
\centering
\begin{tikzpicture}[
    node distance=0.4cm,
    every node/.style={font=\footnotesize, align=center},
    field/.style={draw=accent, thick, minimum width=1cm, minimum height=0.8cm,
                  fill=accentlight}
]
\node[field, minimum width=2.2cm] (ts) {\texttt{timestamp}\\\scriptsize 8\,B};
\node[field, right=of ts, minimum width=2.4cm] (mo) {\texttt{metadata\_offset}\\\scriptsize 8\,B};
\node[field, right=of mo, minimum width=1.4cm] (src) {\texttt{src\_sid}\\\scriptsize 4\,B};
\node[field, right=of src, minimum width=1.4cm] (dst) {\texttt{dst\_sid}\\\scriptsize 4\,B};
\node[field, below=0.6cm of ts.south west, anchor=north west, minimum width=1.6cm] (ia) {\texttt{ingested\_at}\\\scriptsize 4\,B};
\node[field, right=of ia, minimum width=1.0cm] (dt) {\texttt{ds\_tag}\\\scriptsize 2\,B};
\node[field, right=of dt, minimum width=1.0cm] (rt) {\texttt{rel\_tag}\\\scriptsize 1\,B};
\node[field, right=of rt, minimum width=0.9cm, fill=white, dashed] (pad) {pad\\\scriptsize 1\,B};
\end{tikzpicture}
\caption{\texttt{CompactRelation} en \texttt{\#[repr(C)]}: 32~bytes
exactos, alineados a 8. El \emph{layout} se eligió para que un
\texttt{\&[StreamEdge]} (que comparte \emph{layout} byte-a-byte) se
mapee directamente a una región \emph{mmap} cuando el \emph{streaming
backend} \emph{spills} a disco.}
\label{fig:datos-compact}
\end{figure}

\section{Almacenamiento de aristas: \emph{streaming backend} con auto-spill}

El motor mantiene las aristas en un \emph{streaming backend}
per-vertex \emph{degree-adaptive}: \texttt{Inline} para grado $\leq 8$
(en la struct misma, sin alloc), \texttt{Vec<StreamEdge>} para grado
mayor. Cuando el budget global cruza el umbral
(default $10^6$ aristas en \texttt{GraphHunter::new()},
configurable vía \texttt{with\_memory\_budget\_edges}) el backend
escribe un archivo \emph{round-numbered}, lo \emph{mmap}-ea, y
swap-ea cada \texttt{NeighborList} a la variante \texttt{Mmap \{ handle,
byte\_offset, len \}}. El \emph{sort} por \texttt{timestamp} se hace
\emph{durante el write}, así que el archivo es \emph{matcher-ready}.
Las búsquedas binarias \texttt{partition\_point} sobre la \emph{slice}
funcionan indistintamente sobre la región \emph{Inline}, \emph{Vec}
en RAM, o \emph{Mmap} de disco.

\begin{enelcodigo}
\filepath{core/graph-engine/src/streaming.rs}:
\texttt{InMemoryStreamingBackend} con \texttt{append\_edge},
\texttt{finalize}, \texttt{auto\_spill}, \texttt{neighbors\_arc},
\texttt{for\_each\_edge}. Stage 2c part 2 eliminó el antiguo
\texttt{SpillableEdgeStore} (heavy-logs~\#1/\#2) en favor de este
diseño unificado --- el matcher \emph{hot path} y los $\sim$50
\emph{readers} \emph{admin/aggregation} comparten un solo \emph{store}.
\end{enelcodigo}

\section{\texttt{HuntResult} y la migración a \texttt{Arc<str>}}

Un \emph{hunt} produce \texttt{Vec<HuntResult>}, donde
\texttt{HuntResult = Vec<Arc<str>>}. Cada elemento del \texttt{Vec}
es un vértice del \emph{path} resuelto vía el \emph{interner}.

\begin{intuicion}
El cambio de \texttt{Vec<String>} a \texttt{Vec<Arc<str>>} (heavy-logs)
parece menor pero ahorró $\sim 55\%$ del tiempo de \emph{cap-saturated
benches}. \texttt{String::from(...)} \emph{aloca} bytes en \emph{heap};
\texttt{Arc::clone(...)} sólo bumpea un contador atómico. En un
\emph{hunt} con 10K resultados de 4~vértices cada uno, eso son 40K
\emph{allocs} convertidos en 40K \emph{refcount bumps}.
\end{intuicion}

\section{Índices auxiliares}

Sobre el grafo construido el motor mantiene dos índices que aceleran
el \emph{matcher} pero no son la \emph{source of truth}:

\begin{itemize}
    \item \textbf{NLF} (\emph{Neighborhood Label Frequency}): para
    cada vértice, cuántos vecinos salientes tiene de cada
    \texttt{EntityType}. Permite descartar \emph{starts} cuyo
    histograma no satisface el primer paso de la hipótesis.
    \item \textbf{k-hop reachability}: bitmap de los tipos
    alcanzables en a lo sumo \texttt{MAX\_DEPTH=4} saltos, más
    \texttt{min\_first} y \texttt{max\_first} para gating temporal.
\end{itemize}

Ambos se construyen \emph{lazy} en el primer \emph{hunt} sobre la
sesión (heavy-logs~\#3.5) y se invalidan / actualizan
incrementalmente en cada \texttt{add\_relation} (L3 stage~2a).
Capítulo~\ref{cap:indexes} entra en detalle.

\chapter{Pipeline de ingesta}
\label{cap:ingest}

\begin{intuicion}
``Ingesta'' es el camino \texttt{log $\rightarrow$ grafo}. El motor
nunca \emph{procesa} el log entero en RAM: lee batches, parsea cada
batch a \emph{triples} \texttt{(src, rel, dst, ts)}, los empuja al
grafo, repite. Para EVTX de 28~GB esto significa que el peak RAM
queda en el orden del \emph{batch size} (50K registros, $\sim$10~MB)
en lugar del archivo entero.
\end{intuicion}

\section{Arquitectura del \emph{ingest}}

\begin{figure}[H]
\centering
\begin{tikzpicture}[
    node distance=0.7cm and 0.9cm,
    every node/.style={font=\footnotesize, align=center},
    box/.style={draw=accent, thick, rounded corners=2mm,
                minimum height=0.9cm, minimum width=2.4cm,
                fill=accentlight},
    arr/.style={->, >=Latex, accent, thick}
]
\node[box] (src) {\texttt{LogSource}\\\itshape Sentinel, Defender,\\Data Lake, EVTX};
\node[box, right=of src] (poll) {\texttt{SourcePoller}\\\itshape async pump};
\node[box, right=of poll] (parser) {\texttt{LogParser}\\\itshape Sysmon, Sentinel,\\Generic, Forti};
\node[box, right=of parser] (writer) {\texttt{GraphWriter}\\\itshape add\_entity\\add\_relation};
\node[box, right=of writer] (stream) {\texttt{streaming}\\\itshape auto-spill mmap};

\draw[arr] (src) -- (poll);
\draw[arr] (poll) -- (parser);
\draw[arr] (parser) -- (writer);
\draw[arr] (writer) -- (stream);

\node[below=0.3cm of poll, font=\scriptsize\itshape, accent] {watermarks};
\node[below=0.3cm of parser, font=\scriptsize\itshape, accent] {triples};
\node[below=0.3cm of writer, font=\scriptsize\itshape, accent] {sole canonical store};
\end{tikzpicture}
\caption{Pipeline de ingesta. Cada caja es un componente con su
propio módulo en \texttt{platform/sources}, \texttt{platform/parsers},
o \texttt{core/graph-engine}. Las flechas son canales async (\texttt{tokio::sync::mpsc}).}
\label{fig:ingest-pipe}
\end{figure}

\section{De \emph{event} a \texttt{add\_relation}}

\begin{enumerate}
    \item \textbf{Pull batch}. \texttt{SourcePoller} pide un \emph{chunk}
    al \texttt{LogSource} (HTTP, KQL, EVTX file). Devuelve un
    \texttt{RawChunk} con \texttt{rows: Bytes}.
    \item \textbf{Parse}. \texttt{LogParser::parse()} convierte cada
    \emph{row} en cero o más \texttt{ParsedTriple}. Cada \emph{triple}
    es \texttt{(source\_id, source\_type, rel\_type, dest\_id,
    dest\_type, timestamp, metadata)}.
    \item \textbf{Apply al grafo}. \texttt{GraphWriter} llama a
    \texttt{graph.add\_entity} (idempotente — el \emph{interner}
    deduplica) y luego \texttt{graph.add\_relation}. Cada
    \texttt{add\_relation} \emph{single-writes} al
    \texttt{streaming} backend vía \texttt{push\_to\_streaming}
    (post stage~2c part~2 \emph{rip-out}: ya no hay
    \emph{dual-write} a un \texttt{edge\_store} legacy).
\end{enumerate}

\section{Finalización: \texttt{sort\_edges\_by\_timestamp}}

Después de la última \texttt{add\_relation}, el grafo se cierra
explícitamente:

\begin{enelcodigo}
\filepath{core/graph-engine/src/graph.rs}: \texttt{sort\_edges\_by\_timestamp}.
Llama a \texttt{streaming.finalize()} (sortea cada \emph{per-vertex
slice} por \texttt{timestamp}), marca \texttt{edges\_sorted = true}, y
\emph{resetea} las \texttt{OnceLock}s de NLF y \emph{khop\_reach} para
que se construyan \emph{lazy} en el próximo \emph{hunt}.
\end{enelcodigo}

\begin{intuicion}
El \emph{lazy build} de los índices es la decisión que destrabó
el caso 28~GB: antes, \texttt{sort\_edges\_by\_timestamp} construía
NLF y \emph{khop\_reach} \emph{eagerly} y la UI quedaba \emph{frozen}
3~minutos. Ahora termina apenas el \emph{sort per-vertex}; el primer
\emph{hunt} paga el \emph{build}, los siguientes vuelan.
\end{intuicion}

\section{Streaming backend: el \emph{canonical store}}

El \texttt{streaming: Arc<InMemoryStreamingBackend>} es \emph{el} edge
store. Cada \texttt{add\_relation} hace un único \emph{write}:
\texttt{streaming.append\_edge(src, StreamEdge \{ ... \})}, donde
\texttt{StreamEdge} es \emph{layout-compatible} con
\texttt{CompactRelation} (\texttt{\#[repr(C)]}, mismo orden de campos)
para que \emph{readers} aguas-abajo iteren un \texttt{\&[StreamEdge]}
con la misma forma de un \texttt{\&[CompactRelation]} histórico.

El backend es \emph{degree-adaptive} per-vértice (\texttt{Inline} para
grado $\leq 8$, \texttt{Vec<StreamEdge>} para grado mayor) y soporta
\emph{auto-spill incremental}: cuando el \emph{budget global}
(default $10^6$ aristas, configurable) cruza el umbral, el backend
escribe un archivo \emph{round-numbered} \emph{mmap}-eado, sortea cada
slice por \texttt{timestamp} \emph{durante el write}, y swap-ea cada
\texttt{NeighborList} a la variante \texttt{Mmap}. Lecturas siguen
siendo \emph{zero-copy} (\texttt{neighbors\_arc + read guard +
as\_slice}), \emph{matcher partition\_point} sigue siendo \(O(\log d)\).
Detalle en capítulo~\ref{cap:l3-future}.

\section{Incremental NLF update}

El cambio menos visible pero más impactante de stage~2a: cuando una
\texttt{add\_relation} ocurre \emph{después} de que la NLF ya se
construyó, el motor intenta \texttt{NlfTable::try\_record\_edge} en
lugar de invalidar la \texttt{OnceLock}.

\begin{algorithm}[H]
\caption{Incremental NLF update en \texttt{add\_relation}}
\KwIn{\texttt{src\_sid}, \texttt{dst\_sid}, \texttt{dst\_tag}}
\If{\texttt{nlf} está construida}{
    \If{\texttt{nlf.try\_record\_edge(src\_sid, dst\_tag)} == \texttt{true}}{
        \tcc{Saturating add sobre el \emph{heap-backed} counts. NLF
             queda viva.}
    }
    \Else{
        \tcc{Backing está \emph{mmap-spilled} (read-only). Invalidar
             y dejar que se \emph{reconstruya} en el próximo
             \emph{hunt}.}
        \texttt{nlf = OnceLock::new()}
    }
}
\Else{
    \tcc{Nunca se construyó: nada que invalidar.}
}
\end{algorithm}

\begin{intuicion}
Antes de stage~2a, cada \texttt{add\_relation} disparaba un
\emph{rebuild} O($|E|$) de la NLF en el próximo \emph{hunt}. Con
ingestas continuas (Sentinel pulling cada 30~s) eso significa
\emph{rebuild} en cada \emph{hunt} práctico, perdiendo el beneficio
del \emph{lazy build}. Ahora la NLF se mantiene viva mientras el
\emph{backing} sea heap; sólo se invalida cuando es físicamente
imposible mutar (\emph{mmap-spilled}).
\end{intuicion}

\chapter{El \emph{matcher}, tier por tier}
\label{cap:matcher}

\begin{intuicion}
Llamar a \texttt{search\_temporal\_pattern} no ejecuta un solo
algoritmo --- ejecuta el operador físico que el \emph{planner} eligió
para esa hipótesis. Hay cuatro hoy, ordenados por especificidad y
\emph{lift} esperado: \emph{LFTJ-K4} (sólo cuando la hipótesis es un
\(K_4\) homogéneo), \emph{Yannakakis chain}, \emph{Yannakakis tree-walk}
(\emph{shape} \texttt{Graph}), y el \emph{DFS legacy} como
\emph{fallback}. Este capítulo explica cuándo se elige cada uno y por
qué.
\end{intuicion}

\section{El \emph{planner}}

\begin{lstlisting}[language=Rust]
pub enum DispatchPlan {
    DfsMatch,    // default, generic
    K4Lftj,      // M1/S3, sólo K_4 homogéneo
    Yannakakis,  // L1, α-acíclica + flag GRAPHHUNTER_YANNAKAKIS
}

pub fn plan(hypothesis: &Hypothesis) -> DispatchPlan {
    if name_marker_dispatch(hypothesis)   { return K4Lftj; }
    if structural_dispatch(hypothesis)    { return K4Lftj; }
    if yannakakis_dispatch(hypothesis)    { return Yannakakis; }
    DfsMatch
}
\end{lstlisting}

\begin{itemize}
    \item \texttt{name\_marker\_dispatch}: la \emph{back-compat} S3.
    Si el nombre de la hipótesis contiene la cadena ``k4-clique'' y
    \texttt{GRAPHHUNTER\_K4\_LFTJ=1}, ruta al LFTJ.
    \item \texttt{structural\_dispatch} (M1): detecta una hipótesis
    de exactamente 6 pasos homogéneos (mismo \texttt{EntityType},
    misma \texttt{RelationType}, sin predicados) sin requerir la
    palabra mágica.
    \item \texttt{yannakakis\_dispatch} (L1): cualquier hipótesis
    \(\alpha\)-acíclica con \texttt{GRAPHHUNTER\_YANNAKAKIS=1}.
    Ambas \emph{shapes} \texttt{Chain} y \texttt{Graph} aplican.
\end{itemize}

\section{LFTJ-K4: el \emph{worst-case-optimal join} para cliques}

Para una hipótesis \(K_4\) (cuatro vértices, las seis aristas
dirigidas entre todos los pares), el DFS recursivo escala como
\(O(m \cdot \deg^3)\) en el peor caso. \emph{Leapfrog Triejoin}
(Veldhuizen ICDT 2014) se queda en el \emph{AGM bound} \(m^{\rho^*} =
m^2\) — es \emph{worst-case optimal}.

\begin{intuicion}
La idea es construir un \emph{trie} sobre las aristas indexado por
\texttt{(source, dest)} y usar \emph{multi-way intersection} en lugar
de \emph{nested loop}. Para \(K_4\), una vez fijado un par
\((a, b)\), el conjunto de \(c\)s que extiende a un triángulo es
\(N(a) \cap N(b)\) (\emph{intersect SIMD}); el conjunto de \(d\)s que
completa el \(K_4\) es \(\{c\} \cap N(a) \cap N(b) \cap N(c)\).
\end{intuicion}

\begin{enelcodigo}
\filepath{core/graph-engine/src/lftj.rs} y
\filepath{core/graph-engine/src/lftj\_query.rs}. Versiones canónicas
(con \emph{symmetry breaking} hardcoded) en M3++ entregaron el corte
adicional \(|\mathrm{Aut}|\)-veces.
\end{enelcodigo}

\section{DFS legacy con \emph{prunes}}

Cuando ningún operador especializado dispatcha, el motor cae al DFS
recursivo de toda la vida. Pero el DFS de hoy lleva tres \emph{prunes}
acumulados:

\begin{enumerate}
    \item \textbf{NLF}: descarta \emph{starts} cuyo histograma de
    vecinos no satisface el primer paso.
    \item \textbf{k-hop reachability (D5/M2)}: descarta \emph{starts}
    desde los que no puede haber un camino de longitud
    \(\leq \mathrm{MAX\_DEPTH}\) al \texttt{dest\_type} final, dentro
    de la ventana temporal.
    \item \textbf{B9 \emph{symmetry breaking}} para hipótesis
    palindrómicas.
\end{enumerate}

Cada \emph{prune} se aplica antes del DFS principal. Capítulo~\ref{cap:indexes}
detalla cómo se construyen.

\section{Yannakakis chain (L1 stage 2)}

Para cualquier hipótesis \(\alpha\)-acíclica con \emph{shape}
\texttt{Chain}, el motor materializa los \emph{bags} y corre la
\emph{semi-join reducer} de dos pasadas antes de enumerar:

\begin{algorithm}[H]
\caption{Yannakakis chain matcher}
\KwIn{\texttt{bags}: \([B_0, B_1, \ldots, B_{n-1}]\)}
\KwOut{\texttt{Vec<HuntResult>}}
\tcc{Pasada 1 (bottom-up).}
\For{$i$ from $n-2$ \textbf{down to} $0$}{
    \texttt{live\_origins} $\leftarrow \{r.\mathrm{origin} : r \in B_{i+1}\}$\;
    $B_i \leftarrow \{r \in B_i : r.\mathrm{dest} \in \texttt{live\_origins}\}$\;
}
\tcc{Pasada 2 (top-down).}
\For{$i$ from $1$ \textbf{to} $n-1$}{
    \texttt{live\_dests} $\leftarrow \{r.\mathrm{dest} : r \in B_{i-1}\}$\;
    $B_i \leftarrow \{r \in B_i : r.\mathrm{origin} \in \texttt{live\_dests}\}$\;
}
\tcc{Enumerate.}
\For{cada \texttt{row} en $B_0$}{
    \texttt{walk}(\texttt{row.origin}, \texttt{row.dest},
                 \texttt{row.timestamp}, $1$)\;
}
\end{algorithm}

\begin{figure}[H]
\centering
\begin{tikzpicture}[
    node distance=0.3cm,
    every node/.style={font=\scriptsize},
    bag/.style={draw=accent, thick, rounded corners=1mm,
                minimum width=2.4cm, minimum height=1cm,
                fill=accentlight, align=center}
]
\node[bag] (b0) {$B_0$\\\itshape User--Auth--Host};
\node[bag, right=of b0] (b1) {$B_1$\\\itshape Host--Connect--IP};
\node[bag, right=of b1] (b2) {$B_2$\\\itshape IP--Beacon--Domain};

\draw[->, >=Latex, accent, thick, dashed] (b1.south) to[out=-90,in=-90, looseness=1.2] node[below, yshift=-0.2cm]{semi-join} (b0.south);
\draw[->, >=Latex, accent, thick, dashed] (b2.south) to[out=-90,in=-90, looseness=1.2] (b1.south);
\draw[->, >=Latex, accent!60, thick] (b0.north) to[out=90,in=90, looseness=1.2] node[above, yshift=0.2cm]{semi-join} (b1.north);
\draw[->, >=Latex, accent!60, thick] (b1.north) to[out=90,in=90, looseness=1.2] (b2.north);
\end{tikzpicture}
\caption{Las dos pasadas del \emph{semi-join reducer}. Bottom-up
elimina \emph{rows} cuyo \emph{dest} no aparece como \emph{origin} en
el siguiente \emph{bag}; top-down hace lo simétrico. Después de las
dos pasadas, toda \emph{row} sobreviviente participa de al menos un
\emph{full join}.}
\label{fig:matcher-semijoin}
\end{figure}

\section{Yannakakis tree-walk (L1 stage 3b)}

Para \emph{shape} \texttt{Graph} (estrellas, árboles, cualquier
\(\alpha\)-acíclico no-cadena), el \emph{enumerator} camina el
\emph{join tree} construido por GYO. Mantiene un
\texttt{binding: Vec<Option<StrId>>} indexado por \texttt{VarId}; en
cada nodo del árbol consulta el \emph{bag} correspondiente a través
del \emph{shared variable} con el padre.

\begin{algorithm}[H]
\caption{Tree-walk enumerator}
\KwIn{\texttt{join\_tree}, \texttt{bags}, \texttt{binding} (todo \texttt{None})}
\KwOut{\texttt{Vec<HuntResult>}}
\texttt{recurse}(\texttt{node\_idx} = 0)\;
\BlankLine
\Comment{\texttt{recurse(node\_idx)}}
\If{\texttt{node\_idx} == \texttt{join\_tree.nodes.len()}}{
    \texttt{emit}(\texttt{binding} resuelto a \texttt{Vec<Arc<str>>})\;
    \Return
}
\texttt{step} $\leftarrow$ \texttt{join\_tree.nodes[node\_idx]}\;
\texttt{edge} $\leftarrow$ \texttt{hypergraph.edges[step]}\;
\Switch{\texttt{(binding[edge.origin], binding[edge.dest])}}{
    \Case{\texttt{(Some(o), \_)}}{candidates $\leftarrow$ \texttt{by\_origin[step][o]}}
    \Case{\texttt{(None, Some(d))}}{candidates $\leftarrow$ \texttt{by\_dest[step][d]}}
    \Case{\texttt{(None, None)}}{candidates $\leftarrow$ \texttt{bags[step].rows} (root)}
}
\For{\texttt{row} \textbf{in} candidates}{
    intentar \texttt{bind(edge.origin, row.origin)} y \texttt{bind(edge.dest, row.dest)}\;
    \texttt{recurse(node\_idx + 1)}\;
    \texttt{unbind} si correspondía\;
}
\end{algorithm}

\begin{intuicion}
La diferencia clave con el \emph{chain enumerator}: en una cadena
sabemos que \texttt{bag[i+1].origin == bag[i].dest} por la convención
posicional, así que el \emph{lookup} es siempre \emph{by\_origin}. En
una estrella, el \emph{shared variable} con el padre puede aparecer
en cualquiera de los dos lados; por eso el \emph{tree-walk} mantiene
ambos índices y dispatcha por \texttt{(origin\_bound, dest\_bound)}.
\end{intuicion}

\chapter{Índices auxiliares}
\label{cap:indexes}

\begin{intuicion}
NLF y \emph{k-hop reachability} no son la \emph{source of truth} del
grafo --- son resúmenes que el \emph{matcher} consulta para podar
\emph{starts} antes de iniciar el DFS o el \emph{Yannakakis walk}.
Construir el índice cuesta \(O(E)\) (NLF) o \(O(E \cdot \mathrm{MAX\_DEPTH})\)
(\emph{khop\_reach}); el \emph{lazy build} difiere ese costo al
primer \emph{hunt} sobre la sesión.
\end{intuicion}

\section{NLF: \emph{Neighborhood Label Frequency}}

\begin{definition}[NLF]
Para cada vértice \(v\) y \emph{label} \(L \in \Sigma\) (donde
\(\Sigma\) son los \texttt{EntityType}s con \texttt{tag} en
\([0, \texttt{NLF\_TYPE\_COUNT})\)), \(\mathrm{NLF}(v, L)\) cuenta
los vecinos salientes de \(v\) cuyo tipo es \(L\).
\end{definition}

\begin{theorem}[NLF como prune sound]
Sea \(h\) una hipótesis cuyo primer paso requiere \(\geq r_L\)
vecinos de tipo \(L\) en el origen. Si \(\mathrm{NLF}(v, L) < r_L\)
entonces \(v\) no puede ser \emph{start} de ningún \emph{match}.
\end{theorem}

El motor descarta \emph{starts} antes de entrar al DFS aplicando
\(\mathrm{NLF}(v, L) \geq r_L\) para cada \(L\) requerido.

\begin{enelcodigo}
\filepath{core/graph-engine/src/nlf.rs}: \texttt{NlfTable::can\_match}.
El \emph{backing} es \texttt{IndexBacking<[u32; 9]>} --- heap-resident
para grafos chicos, \emph{mmap-spilled} arriba del \emph{budget}
(default 1\,GB).
\end{enelcodigo}

\section{k-hop reachability + extensión temporal (D5/M2)}

NLF cubre 1-hop. Para hipótesis multi-hop el motor también descarta
\emph{starts} desde los que no puede haber un camino de longitud
\(\leq \mathrm{MAX\_DEPTH}\) al \texttt{dest\_type} final.

\begin{definition}[k-hop reachability]
\(\mathrm{bits}[v]\) es un \texttt{u16} donde el \emph{bit} \(L\)-ésimo
está prendido sii existe un camino de longitud
\(\leq \mathrm{MAX\_DEPTH}\) desde \(v\) a algún vértice de tipo \(L\).
\end{definition}

La extensión \emph{M2} agrega dos matrices paralelas
\(\mathrm{min\_first}[v][L]\) y \(\mathrm{max\_first}[v][L]\) que
guardan el \texttt{timestamp} mínimo/máximo del \emph{first edge} de
algún camino válido de \(v\) a \(L\). Permite descartar \emph{starts}
cuya ventana temporal no encaja:

\begin{theorem}[Temporal prune]
Si \(\mathrm{min\_first}[v][L] \geq t_{\mathrm{hi}}\) o
\(\mathrm{max\_first}[v][L] < t_{\mathrm{lo}}\), ningún camino de \(v\)
a \(L\) tiene su primer \emph{edge} dentro de
\([t_{\mathrm{lo}}, t_{\mathrm{hi}})\).
\end{theorem}

\subsection{Storage layout (post heavy-logs \#1+\#2)}

\begin{figure}[H]
\centering
\begin{tikzpicture}[
    node distance=0.4cm,
    every node/.style={font=\footnotesize, align=left},
    box/.style={draw=accent, thick, rounded corners=1mm,
                minimum height=1.2cm, minimum width=3.6cm,
                fill=accentlight}
]
\node[box] (bits) {\texttt{bits: Vec<u16>}\\\scriptsize dense, 2\,B/vertex\\\scriptsize 200\,MB @ 100M vértices};
\node[box, right=of bits] (rank) {\texttt{rank: Vec<u32>}\\\scriptsize dense, 4\,B/vertex\\\scriptsize 400\,MB @ 100M};
\node[box, below=of bits] (mn) {\texttt{min\_first: IndexBacking<[i64;9]>}\\\scriptsize sparse, 72\,B/ranked\\\scriptsize 2.16\,GB @ 30M ranked};
\node[box, right=of mn] (mx) {\texttt{max\_first: IndexBacking<[i64;9]>}\\\scriptsize sparse, 72\,B/ranked\\\scriptsize 2.16\,GB @ 30M ranked};
\end{tikzpicture}
\caption{\texttt{KHopReach} layout. \texttt{rank[v] = NO\_RANK} cuando
\(v\) no tiene aristas salientes (no aporta info temporal); para los
demás, \texttt{rank[v]} indexa filas \emph{sparse} en \texttt{min\_first}
y \texttt{max\_first}. Total: \(\sim 4.9\)\,GB en grafo de 100M
vértices con 30\% \emph{outgoing-having}, vs.\ \(\sim 14.6\)\,GB del
\emph{layout} denso original.}
\label{fig:indexes-khop-layout}
\end{figure}

\subsection{Lazy build (heavy-logs \#3.5)}

Antes de heavy-logs \#3.5, \texttt{sort\_edges\_by\_timestamp}
construía NLF y \emph{khop\_reach} \emph{eagerly}. Para 28~GB EVTX,
ese build duraba 1--3 minutos --- la UI quedaba \emph{frozen}.

Ahora ambos índices viven en \texttt{OnceLock} y se construyen al
\emph{primer hunt} que los consulta:

\begin{lstlisting}[language=Rust]
pub fn ensure_nlf(&self) -> &NlfTable {
    self.nlf.get_or_init(|| NlfTable::build(self))
}

pub fn ensure_khop_reach(&self) -> &KHopReach {
    self.khop_reach.get_or_init(|| KHopReach::build(self))
}
\end{lstlisting}

\begin{intuicion}
La consecuencia visible: abrir una sesión 28~GB es \emph{instantánea}
(sólo \emph{stream} binario + \emph{merge sort}); el primer
\emph{hunt} paga el \emph{build} (\(\sim 30\)\,s NLF + \(\sim 1\)\,min
\emph{khop\_reach}); los siguientes \emph{hunts} vuelan. Si el
analista quiere evitar el primer-hunt-lento, dispara un \emph{hunt}
trivial apenas carga la sesión para \emph{warm-up}.
\end{intuicion}

\section{Incremental NLF (L3 stage 2a)}

\texttt{NlfTable::try\_record\_edge(src\_sid, dst\_tag) -> bool}
bumpea \texttt{counts[src][dst\_tag]} en lugar de invalidar la
\texttt{OnceLock}. Falla sólo cuando el \emph{backing} es
\emph{mmap-spilled} (read-only post-build); en ese caso el
\emph{caller} invalida y deja que el próximo \emph{ensure\_nlf}
reconstruya.

\begin{algorithm}[H]
\caption{\texttt{try\_record\_edge}}
\KwIn{\texttt{src\_sid}, \texttt{dst\_tag}}
\KwOut{\texttt{bool} (true = mutation aplicada o no necesaria)}
\If{\texttt{dst\_tag} \(\geq\) \texttt{NLF\_TYPE\_COUNT}}{\Return \texttt{true}\,(no-op)}
\texttt{slice} $\leftarrow$ \texttt{counts.heap\_slice\_mut()}\;
\If{\texttt{slice} == \texttt{None}}{\Return \texttt{false}\,(mmap)}
\If{\texttt{slice.get\_mut(src\_sid.index())} == \texttt{Some(row)}}{
    \texttt{row[dst\_tag]} $\leftarrow$ \texttt{row[dst\_tag].saturating\_add(1)}\;
}
\Return \texttt{true}\;
\end{algorithm}

\section{Por qué \emph{khop\_reach} delta es más difícil}

NLF es \emph{local}: una arista \((u, v)\) afecta sólo
\(\mathrm{NLF}(u, \cdot)\). \emph{khop\_reach} es \emph{transitiva}:
una arista \((u, v)\) puede cambiar
\(\mathrm{bits}[w]\) para todo \(w\) que tiene un camino a \(u\) en
\(< \mathrm{MAX\_DEPTH}\) saltos.

Capítulo~\ref{cap:l3-future} entra en la solución: la
\emph{delta WCOJ propagation} de Kankanamge-Sahu-Mhedhbi-Salihoglu
SIGMOD 2017.

\chapter{La DSL}
\label{cap:dsl}

\begin{intuicion}
La DSL de Graph Hunter es deliberadamente chica: un analista debe
poder escribir una hipótesis en menos de una línea sin abrir
manual. La gramática base es la cadena con flechas; la única
extensión sintáctica reciente (L1 stage 3c) es el separador
\texttt{;} que abre las puertas a \emph{shapes} no-cadena.
\end{intuicion}

\section{Sintaxis chain (legacy)}

\begin{lstlisting}
User -[Auth]-> Host -[Connect]-> IP {k=2} HAVING count(distinct last) >= 5 WITHIN 30m
\end{lstlisting}

\begin{itemize}
    \item \texttt{Type binding}: cada nodo es un \texttt{EntityType}
    seguido opcionalmente por un identificador (\emph{binding}). Ej.\
    \texttt{User u} declara que el primer vértice es de tipo
    \texttt{User} y se referenciará como \texttt{u} en clauses
    posteriores.
    \item \texttt{-[RelationType\{predicates?\}]->}: arista con
    \emph{relation type} y predicados opcionales sobre la metadata.
    \item \texttt{\{key=value\}} sobre el \emph{dest type}: predicados
    sobre el nodo destino.
    \item \texttt{\{k=N\}} al final: \emph{k-simplicity} (cuántas veces
    puede repetirse un vértice en el \emph{path}).
    \item \texttt{HAVING count(distinct col) >= N}: agregación
    post-\emph{hoc}.
    \item \texttt{WITHIN Ns/Nm/Nh}: ventana temporal entre \emph{first}
    y \emph{last edge} del \emph{path}.
\end{itemize}

\section{Sintaxis Graph (L1 stage 3c)}

\begin{lstlisting}
User c -[Auth]-> Host h1; c -[Auth]-> Host h2; c -[Auth]-> Host h3
\end{lstlisting}

\begin{intuicion}
El primer \texttt{;} marca la hipótesis como \emph{shape}
\texttt{Graph}: cada \texttt{(step, side)} es su propia variable a
menos que un \emph{binding} explícito una posiciones. Una cláusula
después de \texttt{;} puede arrancar de dos formas:

\begin{enumerate}
    \item \emph{Bare binding reference}: un identificador que ya fue
    declarado en una cláusula anterior (\texttt{c} arriba). El motor
    hereda el \texttt{EntityType} del binding original.
    \item \emph{Fresh \texttt{Type binding} declaration}: introduce
    un vértice independiente.
\end{enumerate}
\end{intuicion}

\subsection{Disambiguación}

\begin{itemize}
    \item Identificadores que coinciden con un \texttt{EntityType}
    builtin (\texttt{User}, \texttt{Host}, \texttt{IP}, etc.) siempre
    se parsean como \emph{type declaration}, sombreando bindings del
    mismo nombre.
    \item Identificadores desconocidos producen un error con un
    \emph{hint} listando los bindings declarados.
\end{itemize}

\subsection{Estructura interna del parser}

\begin{figure}[H]
\centering
\begin{tikzpicture}[
    node distance=0.5cm,
    every node/.style={font=\scriptsize, align=center},
    state/.style={draw=accent, thick, rounded corners=1mm,
                  minimum height=0.7cm, minimum width=2.6cm,
                  fill=accentlight}
]
\node[state] (start) {\texttt{parse()}};
\node[state, below=of start] (head) {\texttt{parse\_entity\_type\_with\_binding}};
\node[state, below=of head] (chain) {\texttt{parse\_chain\_steps}};
\node[state, below=of chain] (sep) {peek \texttt{;}};
\node[state, below left=of sep] (clause) {\texttt{parse\_clause\_head}};
\node[state, below right=of sep] (trail) {\texttt{parse\_trailing\_clauses}};

\draw[->, >=Latex, accent, thick] (start) -- (head);
\draw[->, >=Latex, accent, thick] (head) -- (chain);
\draw[->, >=Latex, accent, thick] (chain) -- (sep);
\draw[->, >=Latex, accent, thick] (sep) -- node[left]{si \texttt{;}} (clause);
\draw[->, >=Latex, accent, thick, dashed]
    (clause.east) to[out=0,in=180] (chain.west);
\draw[->, >=Latex, accent, thick] (sep) -- node[right]{else} (trail);
\end{tikzpicture}
\caption{Flujo del parser. \texttt{parse\_chain\_steps} es el bucle
interior reusable; cada \texttt{;} dispara un re-entry vía
\texttt{parse\_clause\_head}, que decide entre \emph{bare binding ref}
y \emph{type declaration}.}
\label{fig:dsl-parser-flow}
\end{figure}

Cómo se lee el flujo en cada nodo del diagrama:

\begin{itemize}
    \item \texttt{parse\_entity\_type\_with\_binding} (\emph{first
    clause head}): consume un \texttt{Type} obligatorio y, si lo
    sigue un identificador no reservado, lo registra como
    \emph{binding}. Es la única forma de declarar el primer vértice
    --- las hipótesis no pueden empezar con un \emph{bare binding ref}
    porque no hay nada declarado todavía.
    \item \texttt{parse\_chain\_steps}: bucle que consume \emph{N}
    pasos \texttt{-[rel]-> Type binding} encadenados. Termina cuando
    ve EOF, \texttt{\{} (suffix \texttt{\{k=N\}} o predicate block),
    \texttt{;} (clause separator), o un identificador
    alfabético (HAVING/WITHIN). Cada terminación delega al manejador
    correspondiente.
    \item \texttt{parse\_clause\_head} (\emph{tras un \texttt{;}}):
    peeka el identificador líder y decide:
    \begin{itemize}
        \item Si matchea un \emph{builtin type name}
        (\texttt{User}, \texttt{Host}, \texttt{IP}, etc.) $\Rightarrow$
        rewind y delega a \texttt{parse\_entity\_type\_with\_binding}
        --- introduce un vértice fresco con su propio binding.
        \item Si matchea un \emph{binding} ya declarado en cláusulas
        anteriores $\Rightarrow$ retorna
        \texttt{(existing\_type, Some(name))} para que el bucle de
        \texttt{parse\_chain\_steps} arranque desde ese vértice
        compartido.
        \item Si no es ninguno de los dos $\Rightarrow$ error con
        \emph{hint} listando los \emph{bindings} válidos. La opción
        de fallback a \texttt{Other(\dots)} se descarta a propósito
        para no swallow-tear errores de tipeo del analista.
    \end{itemize}
    \item \texttt{parse\_trailing\_clauses}: consume \emph{HAVING} y
    \emph{WITHIN} en cualquier orden hasta encontrar \texttt{\{k=N\}}
    o EOF. Estos no pueden re-entrar el bucle de chain steps; son
    estrictamente \emph{post-fix} a la última cláusula del cuerpo.
\end{itemize}

\section{Del AST al matcher}

El AST que el parser produce es un \texttt{Hypothesis}:

\begin{lstlisting}[language=Rust]
pub struct Hypothesis {
    pub name: String,
    pub steps: Vec<HypothesisStep>,
    pub k_simplicity: usize,
    pub aggregations: Vec<AggClause>,
    pub window_seconds: Option<i64>,
    pub shape: HypothesisShape,   // Chain | Graph
}
\end{lstlisting}

El \texttt{shape} es la pieza nueva que destrabó L1 stage 3.
El \emph{matcher} branchea en \texttt{search\_temporal\_pattern\_yannakakis}:

\begin{itemize}
    \item \texttt{Chain}: \texttt{semijoin\_reduce\_chain}
    \(+\) \texttt{enumerate\_chain}.
    \item \texttt{Graph}: \texttt{JoinTree::build}
    \(+\) \texttt{enumerate\_tree}.
\end{itemize}

\section{Plug-and-play en todos los \emph{transports}}

La DSL está expuesta vía \texttt{parse\_dsl} en
\texttt{platform/dsl}. Todos los \emph{transports} la consumen sin
flag adicional:

\begin{itemize}
    \item HTTP: \texttt{POST /run\_hunt} con \texttt{\{ "dsl": "..." \}}.
    \item MCP: \texttt{cmd\_parse\_dsl} +
    \texttt{cmd\_run\_hunt}.
    \item Tauri: \texttt{cmd\_run\_hunt(hypothesis\_json: String)}
    (el frontend serializa el AST).
\end{itemize}

Un cliente que envía una hipótesis con \texttt{;} hoy mismo recibe
una \texttt{Hypothesis::shape = Graph} que el \emph{planner} dispatcha
al \emph{tree-walk} sin tocar nada más.

\chapter{L3 stages 2b y 2c: del diseño al \emph{ship}}
\label{cap:l3-future}

\begin{intuicion}
Este capítulo era el \emph{plano de obra} para 2b (\emph{incremental
khop\_reach}) y 2c (\emph{streaming as source of truth}) cuando se
escribió. Las dos piezas \textbf{ya están entregadas}: stage 2b
\emph{delta WCOJ propagation} corre dentro de \texttt{add\_relation}
y stage 2c hizo el \emph{rip-out} completo del \texttt{edge\_store}
\emph{dual-write}. El \emph{streaming backend} es el único \emph{store};
\texttt{SpillableEdgeStore} fue eliminado del árbol. La narrativa
algorítmica del \emph{plano} sigue siendo la implementación real --- se
mantiene intacta acá como \emph{educational reference}.
\end{intuicion}

\begin{nota}[Estado al cierre de \emph{stage 2c part 2}]
La migración del matcher \emph{hot path} (\texttt{dfs\_match\_smart},
\texttt{dfs\_match\_iterative}, \texttt{edge\_range\_sid}, AMAC pair),
de los $\sim$50 \emph{readers} \emph{admin/aggregation/index-build}
(\texttt{nlf}, \texttt{khop\_reach}, \texttt{lftj}, \texttt{semijoin},
\texttt{yannakakis}, \texttt{anomaly}, \texttt{analytics},
\texttt{gnn\_bridge}, \texttt{matcher-ffi}, \texttt{platform/api}), y el
\emph{rip-out} del \texttt{edge\_store.push} en los cinco \emph{ingest
paths} (\texttt{add\_relation}, \texttt{add\_relations\_batch},
\texttt{replay\_relations}, \texttt{compact\_before},
\texttt{remove\_entities\_and\_relations\_by\_dataset}) entregaron en
conjunto un \(-31\)\% en \texttt{hunt\_latency} contra el \emph{baseline}
\code{pre-2c2}. \texttt{parking\_lot::RwLock} reemplazó a
\code{std::sync::RwLock} en el \emph{streaming backend} para que el par
de \emph{read guards} per-frame quepa en el budget de regresión.
D4 stage 2 (\emph{auto-spill incremental con sort durante write}) se
entregó simultáneamente y queda \emph{plug-and-play} al construir un
\code{GraphHunter} (default \(10^6\) aristas/round). Detalles en
\filepath{docs/perf/latex/cambios-perf-2026-04.tex}.
\end{nota}

\section{Stage 2b — incremental khop\_reach delta-update}

\subsection{El problema}

La \emph{khop\_reach} es transitiva: una arista \((u, v, t)\) cambia
\(\mathrm{bits}[w]\) para todo \(w\) que tiene un camino a \(u\) en
\(< \mathrm{MAX\_DEPTH}\) saltos. La invalidación-y-rebuild que hace
hoy \texttt{add\_relation} cuesta \(O(\mathrm{MAX\_DEPTH} \cdot E)\)
en la próxima \emph{query}; en grafos de 200M aristas, eso son
$\sim$2~min de \emph{rebuild}. Inaceptable para ingestas continuas
(Sentinel pulling cada 30~s).

\subsection{La técnica: Delta WCOJ Propagation}

Kankanamge-Sahu-Mhedhbi-Salihoglu SIGMOD 2017
(``Graphflow: An Active Graph Database'') describe una propagación
incremental que actualiza \(\mathrm{bits}[v]\) y sus
\(\mathrm{min/max\_first}\) en
\(O(\mathrm{MAX\_DEPTH} \cdot d_{\mathrm{reverse}})\) donde
\(d_{\mathrm{reverse}}\) es el \emph{degree} promedio del grafo
inverso. Para grafos \emph{scale-free} típicos esto es
\(\sim 10\)--$100\times$ más barato que un rebuild.

\begin{algorithm}[H]
\caption{\texttt{khop\_reach.delta\_update(u, v, t)} — esquema}
\KwIn{nueva arista \((u, v, t)\), grafo \(G\), índice \texttt{khop\_reach}}
\BlankLine
\Comment{Round 1 — efecto directo: $u$ ahora alcanza $\mathrm{type}(v)$ y todo lo que $v$ alcanza.}
\(\Delta_u \leftarrow \mathrm{bits}[v] \setminus \mathrm{bits}[u]\)\;
\If{\(\Delta_u = \emptyset\) \textbf{and} \(t\) no mejora $\mathrm{min/max\_first}[u]$}{
    \Return\,(la arista no cambia el índice)\;
}
\(\mathrm{bits}[u] \mathrel{|}\!= \mathrm{bits}[v]\)\;
\For{cada \(L\) en \(\mathrm{bits}[v]\)}{
    \(\mathrm{min\_first}[u][L] \leftarrow \min(\mathrm{min\_first}[u][L], t)\)\;
    \(\mathrm{max\_first}[u][L] \leftarrow \max(\mathrm{max\_first}[u][L], t)\)\;
}
\BlankLine
\Comment{Rounds 2..MAX\_DEPTH — propagación reverse.}
\texttt{frontier} \(\leftarrow \{u\}\)\;
\For{$d$ \textbf{from} 2 \textbf{to} \(\mathrm{MAX\_DEPTH}\)}{
    \texttt{next} \(\leftarrow \emptyset\)\;
    \For{\(x\) \textbf{in} \texttt{frontier}}{
        \For{\(w\) \textbf{in} \(\mathrm{predecessors}(x)\) (\emph{reverse adjacency})}{
            \(\Delta_w \leftarrow \mathrm{bits}[x] \setminus \mathrm{bits}[w]\)\;
            \If{\(\Delta_w \neq \emptyset\) \textbf{or} \(t_{wx}\) mejora \(\mathrm{min/max\_first}[w]\)}{
                \(\mathrm{bits}[w] \mathrel{|}\!= \mathrm{bits}[x]\)\;
                actualizar \(\mathrm{min/max\_first}[w][L]\) para cada \(L \in \mathrm{bits}[x]\), usando \(t_{wx}\) (timestamp de la arista \(w \to x\))\;
                \texttt{next} \(\leftarrow\) \texttt{next} \(\cup \{w\}\)\;
            }
        }
    }
    \texttt{frontier} \(\leftarrow\) \texttt{next}\;
}
\end{algorithm}

\subsection{Soundness}

Para que el \emph{prune} siga siendo \emph{sound}, el índice debe ser
una \emph{sobre-aproximación} de la verdadera reachability. La
\emph{delta update} sólo \emph{agrega} bits a \(\mathrm{bits}[w]\) y
sólo \emph{baja} \(\mathrm{min\_first}\) / \emph{sube}
\(\mathrm{max\_first}\), nunca al revés. Un vértice que ya estaba
``alcanza tipo \(L\)'' sigue alcanzándolo --- la invariante se
preserva trivialmente.

\subsection{Costo en mmap-spilled}

Con \texttt{IndexBacking::heap\_slice\_mut} ya implementado, el
\emph{path} \emph{heap-resident} es directo: mutación \texttt{[u][L]}.
Para \emph{mmap-spilled}, el \emph{file} se abre \texttt{MmapMut}
durante el \emph{delta update}, se aplica la mutación, se hace
\texttt{flush()}, y se vuelve a \texttt{Mmap} \emph{read-only}. La
sincronización entre \emph{readers} concurrentes y la mutación
requiere un \texttt{RwLock}; concretamente:

\begin{lstlisting}[language=Rust]
pub struct KHopReach {
    pub bits: Vec<u16>,         // pequeña, siempre heap
    rank: Vec<u32>,             // pequeña, siempre heap
    // Locks per-array para permitir reads concurrentes y mutación
    // exclusiva durante el delta update.
    min_first: RwLock<IndexBacking<[i64; NLF_TYPE_COUNT]>>,
    max_first: RwLock<IndexBacking<[i64; NLF_TYPE_COUNT]>>,
}
\end{lstlisting}

\subsection{Cuándo fallback al rebuild}

Si la cascada \emph{reverse} crece más de un \emph{budget}
(default 100K vértices visitados), el \emph{delta update} aborta y
señala a \texttt{add\_relation} que invalide. Cubre el caso
patológico de un \emph{hub} de millones de predecesores: para esos,
el \emph{rebuild} \(O(E)\) es \emph{cheaper} que la cascada.

\subsection{Tests de equivalencia}

\begin{enumerate}
    \item \emph{Property test}: para un grafo aleatorio + una
    secuencia aleatoria de \texttt{add\_relation}, comparar
    \texttt{KHopReach::build} (rebuild from scratch) contra
    \texttt{delta\_update} aplicado uno a uno. Iguales para todos los
    vértices.
    \item \emph{Soundness test}: ningún \emph{matcher hot path} debe
    podar un \emph{start} válido después de un \emph{delta update}.
\end{enumerate}

\section{Stage 2c — streaming como \emph{canonical store}}

\subsection{El problema}

Hoy \texttt{add\_relation} hace \emph{dual write}: empuja al
\texttt{edge\_store} (canónico) y al \texttt{streaming} (mirror).
Los dos consumen RAM, los dos pagan el costo de cada \emph{insert}.
El plan a largo plazo es eliminar el dual-write: el \emph{streaming
backend} es el almacenamiento canónico, el \texttt{edge\_store} se
construye on-demand sólo cuando un \emph{matcher} legacy lo necesita.

\subsection{Por qué no se hizo en stage 2a}

El \emph{matcher hot path} consulta \texttt{edge\_store} via
\texttt{get\_relations\_by\_sid}, \texttt{iter\_all},
\texttt{neighbors\_of}, \texttt{get\_compact\_relations} --- toda una
familia de \emph{accessors} con semántica \emph{sorted-by-timestamp}
y \emph{binary-search-ready}. Migrar a \texttt{streaming} significa
re-implementar esa familia sobre el nuevo \emph{layout}
\texttt{NeighborList}.

\subsection{El plano}

\begin{enumerate}
    \item \textbf{Adapter layer}. \texttt{streaming::neighbors\_of(sid) ->
    impl Iterator<Item = StreamEdge>} que mimica
    \texttt{get\_compact\_relations}. Sirve como \emph{drop-in
    replacement} con un campo extra \texttt{rel\_type\_tag}.
    \item \textbf{Re-write \emph{matcher hot path}}. Cada uso de
    \texttt{get\_relations\_by\_sid} se reemplaza por la nueva
    función. Los \emph{benches} verifican parity contra la versión
    legacy (target: \(< 10\%\) de regresión).
    \item \textbf{Drop \texttt{edge\_store}}. Con todos los
    \emph{consumers} migrados, \texttt{edge\_store} queda como
    estructura legacy mantenida sólo para
    \texttt{session\_save}/\texttt{session\_load} (donde el formato
    binario \texttt{GHS1} ya escribe \texttt{Relation}s en streaming
    --- la conversión \texttt{NeighborList \(\rightarrow\) Relation}
    al guardar es directa).
\end{enumerate}

\subsection{Beneficios esperados}

\begin{itemize}
    \item \textbf{RAM}: una sola copia de las aristas en lugar de dos.
    Para 28~GB EVTX \(\sim\) 100M aristas \(\times\) 32~B \(\approx\)
    3.2~GB ahorrado.
    \item \textbf{Latencia de \emph{insert}}: una sola escritura por
    \texttt{add\_relation} (en lugar de dos).
    \item \textbf{Concurrencia}: el \emph{streaming backend} ya
    soporta \emph{readers} concurrentes vía
    \texttt{RwLock<HashMap<\ldots>>}; los \emph{matchers} de varios
    \emph{transports} pueden correr simultáneamente.
\end{itemize}

\subsection{Riesgos y mitigaciones}

\begin{itemize}
    \item \textbf{Regresión en \emph{matcher hot path}}: el
    \texttt{edge\_store} usa \emph{contiguous Vec} sortado, perfecto
    para \emph{cache prefetch}; la \texttt{NeighborList} tiene tres
    variantes (\texttt{Inline}/\texttt{Vec}/\texttt{Btree}) con \emph{cache
    behavior} mixto. Mitigación: benches contra \texttt{matching\_hops}
    y \texttt{hunt\_latency} antes de cada commit; si la regresión
    \(> 10\%\), revisar el \emph{layout} de la \emph{Vec variant}.
    \item \textbf{Compatibilidad de \emph{snapshot}/\emph{load}}: el
    formato \texttt{GHS1} actual lee \texttt{Relation}s en
    \emph{streaming} y los inserta vía \texttt{add\_relation}. Eso
    sigue funcionando porque \texttt{add\_relation} popula el
    \emph{streaming backend} sin más; el \texttt{edge\_store} legacy
    queda vacío en \emph{post-load} y los \emph{matchers} legacy
    paulatinamente migran a \emph{streaming}.
\end{itemize}

\section{Cronograma estimado}

\begin{itemize}
    \item Stage 2b (delta khop\_reach): \(\sim 2\) semanas. La
    propagación inversa es bien-comprendida; el grueso es el
    \emph{property test} de equivalencia delta-vs-rebuild.
    \item Stage 2c (canonical streaming): \(\sim 3\) semanas. La
    re-escritura del \emph{matcher hot path} toca 12+ \emph{call
    sites} en \filepath{core/graph-engine/src/graph.rs} más los
    \emph{benches}; el riesgo de regresión es lo que más empuja el
    timeline.
\end{itemize}

\chapter{Frontend y plug-and-play}
\label{cap:plug-and-play}

\begin{intuicion}
La regla \emph{plug-and-play} dice: cada optimización S/M/L se
activa por defecto al levantar el binario. Sin \emph{flags} a setear,
sin \emph{config} a tocar, sin \emph{toggles} en la UI. El usuario
levanta \texttt{npm run tauri dev} y todo lo que está \emph{shipped}
se nota inmediatamente. Este capítulo documenta cómo se materializa
esa promesa.
\end{intuicion}

\section{Flags default-on al \emph{startup}}

\filepath{apps/tauri/src-tauri/src/lib.rs} y
\filepath{apps/cli/src/main.rs} corren un loop al inicio que setea
flags si el usuario no las exportó:

\begin{lstlisting}[language=Rust]
for (key, default) in [
    ("GRAPHHUNTER_LFTJ_PLANNER", "1"),  // M1 structural K_4
    ("GRAPHHUNTER_K4_LFTJ", "1"),        // S3 name marker
    ("GRAPHHUNTER_YANNAKAKIS", "1"),     // L1 stage 2/3
] {
    if std::env::var(key).is_err() {
        unsafe { std::env::set_var(key, default) };
    }
}
\end{lstlisting}

Override con \texttt{KEY=0} sigue funcionando para usuarios que
quieran el \emph{path} legacy (debugging, A/B comparison).

\section{La pildora del toolbar}

\begin{figure}[H]
\centering
\begin{tikzpicture}[
    node distance=0.4cm,
    every node/.style={font=\scriptsize},
    pill/.style={draw, rounded corners=4mm, minimum width=2.0cm,
                 minimum height=0.7cm, font=\footnotesize}
]
\node[pill, draw=accent, fill=accentlight] (perf) {Perf: M+L1};
\node[pill, draw=green!50!black, fill=green!10, right=of perf]
    (memok) {Mem: 14\%};
\node[pill, draw=orange!70!black, fill=orange!10, right=of memok]
    (memwarn) {Mem: 78\%};
\node[pill, draw=red!70!black, fill=red!10, right=of memwarn]
    (memcrit) {Mem: 92\%};
\end{tikzpicture}
\caption{Pildoras del \emph{toolbar}. La de \emph{perf} indica si
los \emph{flags} M+L1 están activos (verde) o si el usuario forzó el
\emph{legacy path} (gris). La de \emph{mem} cambia de verde a ámbar
(\(\geq 75\%\) RSS/total) a rojo (\(\geq 90\%\)).}
\label{fig:plug-pills}
\end{figure}

\subsection{Implementación}

\filepath{apps/tauri/src-tauri/src/commands/perf.rs} expone
\texttt{cmd\_get\_perf\_status} que retorna una struct con todos los
flags resueltos. \filepath{apps/tauri/src-tauri/src/commands/memory.rs}
expone \texttt{cmd\_get\_memory\_status} con RSS + host total +
ratio.

\filepath{apps/tauri/src/App.tsx} llama ambos comandos al \emph{mount}
y polea el de memoria cada 5~s. Las pildoras del \emph{toolbar} se
renderizan con \emph{data-attributes} que \texttt{App.css} mapea a
los tres estados.

\section{Telemetría adicional en \emph{batch}}

El endpoint \texttt{run\_hunt\_batch} (L2) retorna una respuesta
con tres \emph{cache counters} que el \emph{frontend} puede mostrar
como \emph{tooltip} sobre la pildora de \emph{perf}:

\begin{lstlisting}[language=Rust]
pub struct RunHuntBatchResponse {
    pub entries: Vec<RunHuntBatchEntry>,
    pub indexes_were_warm: bool,
    pub bag_cache_hits: usize,         // L2 stage 2
    pub bag_cache_misses: usize,
    pub result_cache_hits: usize,      // L2 stage 3a
    pub last_prefix_resume_depth: usize,  // L2 stage 3b
    pub prefix_cache_hits: usize,
}
\end{lstlisting}

\section{Doc: ¿qué tocó el usuario?}

La promesa \emph{plug-and-play} se hace verificable: si después de
una optimización el usuario tuvo que cambiar algo (env var, config
file, \emph{toggle}), la promesa se rompió. La regla quedó codificada
en el \emph{memory feedback} \filepath{feedback\_perf\_plug\_and\_play.md}:

\begin{enumerate}
    \item \textbf{Default-on al \emph{startup}}: cada \emph{flag}
    nuevo se setea en \texttt{lib.rs} y \texttt{main.rs} con
    \texttt{if std::env::var(key).is\_err()}.
    \item \textbf{Surface en UI}: extender
    \texttt{cmd\_get\_perf\_status} y la pildora del toolbar.
    \item \textbf{Default-off sólo con causa}: \emph{flags} que
    reescriben el DFS \emph{hot loop} (AMAC, \emph{semi-join})
    quedan opt-in durante un \emph{soak window}; cuando los
    \emph{parity tests} aguantan \(\geq 1\) semana, se flipean.
    \item \textbf{Doc obliga}: cada \emph{flip} actualiza
    \filepath{08-env-vars.tex}, \filepath{12-preguntas-frecuentes.tex}
    y \texttt{make all} para refrescar los 4 PDFs.
    \item \textbf{Capabilities de \emph{build}}: \emph{Cargo features}
    como \texttt{symmetry-nauty} se reflejan via constantes
    \texttt{HAS\_NAUTY}-style en \texttt{cmd\_get\_perf\_status}.
\end{enumerate}

\section{El cierre del libro}

Esta es la última pieza pedagógica del diseño. El resto de los libros
GraphHunter (\emph{Arquitectura}, \emph{Especificación}, \emph{Manual},
\emph{Paper}) profundizan en aspectos específicos:

\begin{itemize}
    \item \emph{Arquitectura}: el \emph{stack} de \emph{crates}, las
    capas, las decisiones de \emph{layering}.
    \item \emph{Especificación}: los \emph{contracts} de HTTP, MCP,
    FFI, las variables de entorno, los formatos de \emph{log} aceptados.
    \item \emph{Manual}: el \emph{user manual} para el analista ---
    cómo se ingiere un dataset, cómo se escribe una hipótesis, cómo se
    interpretan los resultados.
    \item \emph{Paper}: la justificación académica con
    \emph{benchmarks}, \emph{property tests}, y comparaciones contra
    el estado del arte.
\end{itemize}

Este libro es la guía de diseño que los une.


\end{document}
